\begin{corollary}\label{cor:x19smooth}
The full semiuniversal deformation space of a $\Delta_{19}$-regular
anticanonical hypersurface $X_{19}$ is smooth of dimension $30$.
The threefold $X_{19}$ has a convergent smoothing.
\end{corollary}

\begin{proof}
We give the finite relation calculation explicitly.  Order its nodal
faces by the following vertex sets, using \eqref{eq:V19}:
\[
\begin{array}{rrrr@{\qquad}rrrr@{\qquad}rrrr}
0&1&3&5&0&2&4&5&6&8&9&15\\
7&10&11&16&1&6&13&18&1&7&13&17\\
2&6&12&18&2&7&12&17&3&9&10&13\\
4&8&11&12&3&9&14&15&3&10&14&16\\
4&8&14&15&4&11&14&16&&&&
\end{array}
\]
For each set take the primitive relation on its four rays whose
first nonzero coefficient is positive, and extend it by zero on the
other rays.  Let $N$ be the resulting $14\times19$ matrix.  The
degree-seven root has divisor row
\[
 r=(0,0,0,0,0,0,0,0,0,0,0,0,0,0,0,1,-1,1,-1).
\]
Exact row reduction gives
\[
 \rk N=12,\qquad \rk\begin{pmatrix}N\\r\end{pmatrix}=12.
\]
A relation in the latter matrix, nonzero in all fifteen coordinates, is
\[
 (-2,2,-1,1,1,-3,-2,4,-2,2,3,1,-1,-3,-1).
\]
Replace $r$ by a basis of the full two-dimensional relation space
on the pentagon rays $v_{14},\ldots,v_{18}$.  The resulting
$16\times19$ matrix has rank $13$.  Thus both the selected and the
full local relation kernels have dimension three; they are equal.
These matrices are reconstructed from the printed vertices and face
sets in the exact computation listed in Appendix~\ref{app:cert}.

The rows detect homology relations.  Indeed, the correction term for
non-toric divisor classes on a smooth crepant resolution vanishes:
every singular two-face has a dual edge of lattice length one and
every other two-face is a unimodular triangle.  Toric divisors
therefore span $H^2$ by Batyrev's formula \cite{Batyrev94}.  The
inserted divisor over the pentagon pairs to zero with $K_E^\perp$.
Consequently the full link kernel $K$ is already $K_S$, and the
displayed vector proves smoothing by Theorem~\ref{thm04:criterion}.

For the dimension, the polar has $32$ lattice points, no
facet-interior lattice points, and zero correction term in Batyrev's
$h^{2,1}$ formula, since the original edges are primitive.  Hence
$h^{2,1}(Y_0)=32-5=27$.  One also has
\[
 f_{0*}\Theta_{Y_0}=\Theta_{X_{19}},\qquad
 R^1f_{0*}\Theta_{Y_0}=0.
\]
Vector fields lift to the vertex blow-up of a cone and to a chosen
nodal small resolution.  For the higher direct image, formal
functions reduces the degree-seven assertion to
\[
 H^1(E,\Theta_E\otimes\cO_E(-kK_E))=
 H^1(E,\cO_E((1-k)K_E))=0\qquad(k\ge0).
\]
The first equality follows from
$\Theta_E\cong\Omega^1_E\otimes\cO_E(-K_E)$ and toric Bott
vanishing.  The second follows from rationality and Serre duality
for $k=0,1$, and Kodaira vanishing for $k\ge2$.  At a node the
graded terms are copies of $\cO(k+2)$ and $\cO(k-1)$ on $\PP^1$,
whose first cohomology vanishes.  Leray and contraction with the
volume form give $h^1(X_{19},\Theta_{X_{19}})=27$.
Theorem~\ref{thm:D} now gives $\dim\mathbb T^1_{X_{19}}=27+3=30$.

The construction in the proof of Theorem~\ref{thm04:smoothbase}
gives a smooth formal source whose blow-down differential covers
$V_S$.  Here $V_S=\mathbb T^1_{X_{19}}$ because $K_S=K$.
The same independent-coordinate argument therefore makes the
\emph{full} deformation ring regular, not only its selected
quotient.  Regularity of the completed analytic ring implies
smoothness of the analytic germ, with the dimension just computed.
\end{proof}

Theorem~\ref{thm:A} and Corollary~\ref{cor:x19smooth} thus concern
different deformation constructions: $X_{19}$ is smoothable, while
every single-degree ambient family in the stated class retains a
singular anticanonical fibre.
